\documentclass[10pt,a4paper]{amsart}
\usepackage[margin=19mm]{geometry}
\usepackage{amsmath,amssymb,mathtools}
\usepackage[scr=rsfs,cal=euler]{mathalfa}
\usepackage{xfrac}
\usepackage[unicode,hidelinks,bookmarks=false]{hyperref}
\newcommand{\ie}{i.e.\ }
\newcommand{\cf}{cf.\ }
\newcommand{\resp}{respectively\ }
\newcommand{\Subsection}{\subsection}
\providecommand{\nobreakdash}{\nobreak\hbox{-}}
\setlength{\emergencystretch}{2em}
\multlinegap0pt
\usepackage{tikz}
\usepackage{amssymb}
\usepackage[normalem]{ulem}
\newtheorem{lem}{Lemma}
\newtheorem{thm}{Theorem}
\newtheorem{rem}{Remark}
\title{On the fundamental spectral gap of weighted Schrodinger operators}
\author{Mohammed Ahrami}
\date{Source: arXiv:2605.24560v1 (CC BY 4.0)}
\begin{document}
\maketitle

\noindent\emph{Source and licence.} On the fundamental spectral gap of weighted Schrodinger operators. Version arXiv:2605.24560v1 (CC BY 4.0), distributed by arXiv under the Creative Commons Attribution 4.0 International licence, http://creativecommons.org/licenses/by/4.0/, as recorded in the arXiv licence metadata for this exact version and shown on the abstract page https://arxiv.org/abs/2605.24560v1. Full licence notice: \texttt{licenses/arxiv-2605.24560-CC-BY-4.0.txt}.\par\smallskip

\noindent\textit{Editorial scope.} Counted unit: the article's thm step2 (source0.tex lines 580--602). The article's complete original proof of it is delivered with it (source0.tex lines 603--729, one \texttt{proof} environment of 127 lines). No mathematics is written, repaired or re-derived: the statement and the proof are the article's own lines, in the article's own notation, and no retained line is substituted or re-flowed.  Retained uncounted as the source-local context the proof needs: lines 115--118; lines 241--271; lines 278--286; lines 357--360; lines 419--429; lines 571--578.  Nothing was omitted from any retained span, so the package carries no omission record.  No retained line contains a citation, so the wrapper carries no bibliography and no bibliography entry.  No retained line contains a figure environment, an image-inclusion command or drawing code, so no image is delivered and the package declares no asset dependency.  Editorial formatting only, in addition: an AMS \texttt{amsart} preamble that restates the article's own declarations and macros used by the retained lines, namely the environments \texttt{lem}, \texttt{thm}, \texttt{rem} on separate counters, together with the standard packages those lines need.  The retained environments print in this document as Lemma 1, Theorem 1, Remark 1 in order of appearance, whereas the article prints the same items with the numbering of its own full document.  The complete native source of the article as distributed in the arXiv e-print for this exact version is bundled unchanged as \texttt{sources/201284/source0.tex}.
\begin{equation}\label{ePnoP}
H(V,w) u :=  -u^{\prime \prime} +V(x) u  = \lambda  w(x) u,
\quad\quad x\in [0,\pi].
\end{equation}

\begin{lem}\label{F-H}
Suppose that $V(.,\kappa)$ and $w(.,\kappa)$ are one-parameter families of real-valued, locally $L^1$ functions, \( C^1 \) in \( \kappa \) for \( \kappa \) in some real interval, with
$\inf_{x} V(x,\kappa) > - \infty$, $C \ge w(x,\kappa) \ge \frac 1 C$ for some $C > 0$, and
$\frac{\partial V}{\partial \kappa}(x,\kappa)$
and $\frac{\partial w}{\partial \kappa}(x,\kappa) \in L^1(0, \pi)$. Without loss of generality we standardize the first two normalized eigenfunctions $u_{1,2}(x)$ associated with 
$\lambda_{1,2}$
so that
$u_{1,2}(x) > 0$ for $0 < x < \epsilon$ for some $\epsilon$.  Then
\begin{enumerate}
\item \[
\frac{d\lambda_{n}(\kappa)}{d\kappa}=-\lambda_{n}\int_{0}^{\pi}\frac{\partial w}{\partial \kappa}(x,\kappa)u_{n}^{2}(x,\kappa)dx+\int_{0}^{\pi}\frac{\partial V}{\partial \kappa}(x,\kappa)u_{n}^{2}(x,\kappa)dx .
\]
\item
$\dfrac{u_{2}}{u_{1}} $ is decreasing on $(0,\pi)$.

\item The equation $ \vert u_{1}(x)\vert=\vert u_{2}(x)\vert $ has
either one or two solutions on $(0,\pi)$.
 \item There exist two points $x_{-}$ and $x_{+}$,
  $ 0 \leq x_{-}< x_{+} \leq \pi $, at least one of which is interior to $(0, \pi)$, such that 
$u_{1}^{2}(x)>u_{2}^{2}(x)$ on $(x_{-},x_{+})$ and $u_{1}^{2}(x)\le u_{2}^{2}(x)$ on $(x_{-},x_{+})^c$.
\item
The equation $ \lambda_1 \vert u_{1}^2(x)\vert=\lambda_2 \vert u_{2}^2(x)\vert $ has
either one or two
solutions on $(0,\pi)$.
 \item 
There exist two points $\widehat{x}_{-}$ and $\widehat{x}_{+}$,
  $ 0 \leq \widehat{x}_{-}< \widehat{x}_{+} \leq \pi $, at least one of which is interior to $(0, \pi)$, such that 
$\lambda_1 u_{1}^{2}(x)> \lambda_2 u_{2}^{2}(x)$ on $(\widehat{x}_{-},\widehat{x}_{+})$ and $\lambda_1 u_{1}^{2}(x)\le \lambda_2 u_{2}^{2}(x)$ on $(\widehat{x}_{-},\widehat{x}_{+})^c$.
\end{enumerate}

\end{lem}

\begin{lem}\label{G'''}
Let $G:[0,\pi]\to\mathbb{R}$ be a function such that $G^{''}(x)$ is absolutely continuous. If $u$ satisfies \eqref{ePnoP}, where $V(x),w(x)$ are differentiable, then
\begin{align*}
G(\pi)[u'^{2}(\pi)&+(\lambda w(\pi)-V(\pi))u^{2}(\pi)]-G(0)[u'^{2}(0)+(\lambda w(0)-V(0))u^{2}(0)]\\
&\quad\quad\quad\quad+ \dfrac{1}{2}[G''(\pi)u^{2}(\pi)-G''(0)u^{2}(0)]\\
&=\int_{0}^{\pi}\left[2G'(\lambda w-V)+G(\lambda w'-V')+\frac{G'''}{2}\right]
u^{2}(x)dx.
\end{align*}
\end{lem}

\begin{align}\label{s-l}
    {-u''+V(x)u=\lambda w(x)u}&\\
    u(0)=u(\pi)=0&,\nonumber
\end{align}                                                   

\begin{thm}\label{step}
For a given bounded measurable, strictly positive weight function $w$, and for every convex 
and non-affine potential $V$ (i.e. not of the form $ax+b$), there exists a linear potential $V_{a}=ax$ such that
$$
 \Gamma[V,w] \geq  \Gamma[V_{a},w].$$
Alternatively, for a given  
bounded measurable potential function $V$, and for every concave and not-affine
weight $w$ (i.e. not of the form $mx+p$), there exists an affine weight $w_m(x)=mx+p$ such that 
$$
 \Gamma[V,w] \geq  \Gamma[V,w_m].$$  
\end{thm}

\begin{rem}\label{rem12}
    In particular, for the affine weight $w(x)=mx+p$, $m>0$, in the interval $[0,\pi]$, the function $w$ increases since $w^{\prime}(x)=m>0$.
Therefore $\min_{x} w(x)=w(0)=p.$ Then

\[
\int_0^{\pi}x^2(u_2^2(x) -u^2_1(x))dx\leq\frac{2\pi^2}{p}.
\]
\end{rem}

\begin{thm}\label{step2}
Let $w(x)=mx+p$ be a fixed positive affine weight function with $m>0$. Assume that
\[
p>5m\pi
%\int_0^{\pi}x^2(u_2^2(x) -u^2_1(x))dx<\frac{2\pi}{5m}.
\]

\noindent Then any minimizer of the fundamental gap
\[
\Gamma[V,w] = \lambda_2(V,w) - \lambda_1(V,w)
\]
over the class of convex potentials $ \mathcal{C}_M$ must be constant.


%For affine potentials $V(x)=ax$ and for affine positive weights $w(x)=mx+p$ with $m\ne0$, and we assume that $\Gamma[V,w]<\frac{a}{m}$, then any 
%gap-minimizing
%optimal potential $V_{*} \in \mathcal{C}_M  $ is constant.
 % Alternatively, for a given  
%bounded measurable potential function $V$, the 
%gap-minimizing
%optimal weight $w_*\in \mathcal{S}_{N_{<,>}}$ is constant $\text{if } \lambda< \frac{t}{m}$.
    
\end{thm}

\begin{proof}  
By Theorem \ref{step}, for any convex potential $V$, there exists an affine potential
\[
V_a(x) = ax 
\]
such that
\[
\Gamma[V,w] \geq \Gamma[V_a,w].
\]

%Thus, it suffices to consider affine potentials.

\noindent Let the first two normalized eigenfunctions $u_{1,2}(x)$ associated with the first eigenvalues $\lambda_{1,2}$ of problem (\ref{s-l}). 

\noindent We consider a family of potentials $V^{\theta}=\theta ax$. 
 From Lemma \ref{F-H}, we have
 

\begin{align*}
&\frac{d(\lambda_2(V^{\theta})-\lambda_1(V^{\theta}))}{d\theta}\\
&=\int_0^{\pi} \frac{\partial V^{\theta}}{\partial \theta}[u_2^2(x)-u_1^2(x)]dx\\
 &=a\int_0^{\pi}x[u_2^2(x)-u_1^2(x)]dx.
\end{align*}
 

\noindent Suppose that a critical point of the gap occurs
at some $a>0$, then
$$
\int_0^{\pi}x(u_2^2(x)-u_1^2(x))dx =0.
$$
\noindent Applying Lemma \ref{G'''} for $G(x)=x$, we have
\begin{align*}
&\pi[(u_n^{\prime}(\pi))^{2}+(\lambda_n w(\pi)-a\pi )u_n^{2}(\pi)]\\
 &=\int_0^\pi [2(\lambda_n w-ax)+x(\lambda_n w^{\prime}-a)]u_n^2
(x)dx.
\end{align*}

\noindent Since $u_n$ satisfies the Dirichlet boundary conditions $u_n(0)=u_n(\pi)=0$, 

$$\pi(u_n^{\prime}(\pi))^{2}
 =\int_0^\pi [2(\lambda_n w-ax)+x(\lambda_n w^{\prime}-a)]u_n^2
(x)dx.
$$ 

\noindent Similarly for $G(x)=x^2$ and from the Dirichlet boundary conditions, we find that
$$\pi^2(u_n^{\prime}(\pi))^{2}\\
=\int_0^\pi [4x(\lambda_n w-ax)+x^2(\lambda_n w^{\prime}-a)]u_n^2(x)dx.$$
\noindent Consequently
\begin{align*}
&\int_0^\pi [2(\lambda_1 w-ax)+x(\lambda_1 w^{\prime}-a)]u_1^2
(x)dx\\
&=\frac{1}{\pi}\int_0^\pi [4x(\lambda_1 w-ax)+x^2(\lambda_1 w^{\prime}-a)]u_1^2(x)dx
\end{align*}
\noindent and
\begin{align*}
&\int_0^\pi [2(\lambda_2 w-ax)+x(\lambda_2 w^{\prime}-a)]u_2^2
(x)dx\\
&=\frac{1}{\pi}\int_0^\pi [4x(\lambda_2 w-ax)+x^2(\lambda_2 w^{\prime}-a)]u_2^2(x)dx
\end{align*}





%\noindent By theorem \ref{step}, we have $$
 %\Gamma[V,w] \geq  \Gamma[ax,mx+p].$$
\noindent  Subtracting the two identities corresponding to $n=1$ and $n=2$, and using $
\int_0^{\pi}x(u_2^2(x)-u_1^2(x))dx =0$, we obtain

%\noindent We note that $\int_0^\pi w(x)u_n^2(x)dx=1$, \noindent since $
%\int_0^{\pi}x(u_2^2(x)-u_1^2(x))dx =0$.
%\noindent Combining the previous identities
\begin{align*}
&  
\frac{2\pi}{5}(\lambda_2-\lambda_1)\\
&=\left(m(\lambda_2-\lambda_1) -a\right)\int_0^{\pi}x^2(u_2^2(x) -u^2_1(x))dx.\\
\end{align*}

\noindent Hence
\begin{align*}
&  
(\lambda_2-\lambda_1)  \left(\frac{2\pi}{5}-m\int_0^{\pi}x^2(u_2^2(x) -u^2_1(x))dx\right)\\
&=-a\int_0^{\pi}x^2(u_2^2(x) -u^2_1(x))dx.\\
\end{align*}


\noindent Let
\[
\phi(x) = x^2 - Ax - B
\]
\noindent chosen such that
\[
\phi(x_-) = \phi(x_+) = 0.
\]

\noindent Then


$$\displaystyle  \left\{%\left( 
    \begin{array}{ll}
   
   \phi(x) \leqslant  0 \qquad\text{on}~~(x_-,x_+)^c,
   \\
   \phi(x)\geqslant 0 \qquad\text{on}~~(x_-,x_+).                                                                                      \end{array}                                                                                \right.$$

\noindent We obtain
\[
\int_0^\pi x^2 (u_2^2 - u_1^2)\,dx
= \int_0^\pi \phi(x)(u_2^2 - u_1^2)\,dx
+ A \int_0^\pi x(u_2^2 - u_1^2)\,dx
+ B \int_0^\pi (u_2^2 - u_1^2)\,dx>0.
\]

\noindent By Remark \ref{rem12} \[
\int_0^{\pi}x^2(u_2^2(x) -u^2_1(x))dx\leq\frac{2\pi^2}{p}.
\]
It follows that \[\frac{2\pi}{5m}
\leq 
 \frac{2\pi^2}{p}.
\]
But the assumption
\[
p>5m\pi
\]
which yields a contradiction, since the left-hand side is positive. Then the gap-minimizing
optimal potential $V_{*} \in \mathcal{C}_M $ is constant.
\end{proof}  

\end{document}
